\documentclass[10pt]{article}
\usepackage[T1]{fontenc}
\usepackage[utf8]{inputenc}
\usepackage[french]{babel}
\usepackage{amsmath,amssymb,mathrsfs,mathpazo}
\usepackage{geometry}
\geometry{a4paper,margin=2cm}
\usepackage{graphicx,array,tabularx,tabularray,colortbl,multirow,multicol,booktabs,enumitem,pifont,ulem,fancybox,diagbox,cancel}
\usepackage{pstricks,pstricks-add,pst-plot,pst-tree,pst-node,pst-text,pst-func,pst-math,pst-bspline,pst-3dplot,pst-eucl,pst-solides3d}
\let\pscalculate\undefined
\usepackage{pst-bezier}
\usepackage{tikz,pgfplots,tkz-tab}
\usetikzlibrary{arrows,arrows.meta,patterns,calc,positioning,trees,decorations.markings,shapes.geometric,shapes.misc,shapes.arrows,intersections,decorations.pathmorphing,decorations.pathreplacing,backgrounds,shadows,babel}
\pgfplotsset{compat=1.18}
\usepgfplotslibrary{fillbetween}
\usepackage[np]{numprint}
\usepackage{listings,makecell,pas-tableur,twoopt}
\let\barre\undefined
\usepackage[pstricks]{tablvar}
\definecolor{darkgreen}{RGB}{0,100,0}
\providecommand{\red}{\color{red}}
\providecommand{\blue}{\color{blue}}
\providecommand{\green}{\color{green}}
\providecommand{\black}{\color{black}}
\providecommand{\white}{\color{white}}
\providecommand{\gray}{\color{gray}}
\providecommand{\cyan}{\color{cyan}}
\providecommand{\magenta}{\color{magenta}}
\providecommand{\yellow}{\color{yellow}}
\providecommand{\violet}{\color{violet}}
\providecommand{\orange}{\color{orange}}
\providecommand{\euro}{\texteuro}
\providecommand{\EUR}{\texteuro}
\providecommand{\ieme}{\textsuperscript{e}}
\providecommand{\degres}{\ensuremath{^\circ}}
\providecommand{\degree}{\ensuremath{^\circ}}
\providecommand{\vv}[1]{\overrightarrow{#1}}
\providecommand{\Coord}[1]{#1}
\providecommand{\decosix}{}
\providecommand{\starredbullet}{\ensuremath{\bullet}}
\pagestyle{empty}
\setlength{\parindent}{0pt}
\usepackage{adjustbox}\DeclareUnicodeCharacter{2236}{:}\DeclareUnicodeCharacter{2190}{\ensuremath{\leftarrow}}
\providecommand{\R}{\mathbb{R}}
\providecommand{\N}{\mathbb{N}}
\providecommand{\Z}{\mathbb{Z}}
\providecommand{\Q}{\mathbb{Q}}
\providecommand{\D}{\mathbb{D}}
\providecommand{\C}{\mathbb{C}}
\providecommand{\e}{{\mathrm{e}}}
\providecommand{\vect}[1]{\overrightarrow{#1}}
\providecommand{\vectt}[1]{\overrightarrow{\mathrm{#1}}}
\providecommand{\vvt}[1]{\overrightarrow{\mathrm{#1}}}
\providecommand{\np}[2][]{{#2~#1}}
\providecommand{\vv}[1]{{\overrightarrow{#1}}}
\renewcommand{\d}{{\mathrm{d}}}
\providecommand{\degres}{{^\circ}}
\providecommand{\degre}{{^\circ}}
\providecommand{\blue}{\color{blue}}
\providecommand{\red}{\color{red}}
\providecommand{\enskip}{\quad}
\providecommand{\llbracket}{\lbrack\!\lbrack}
\providecommand{\rrbracket}{\rbrack\!\rbrack}
\providecommand{\euro}{{\text{€}}}
\providecommand{\no}{{\text{n°}}}
\providecommand{\interval}[2]{{[#1;#2]}}
\providecommand{\intervFO}[2]{{[#1;#2[}}
\providecommand{\dotfill}{\ldots}

\providecommand{\ds}{\displaystyle}
\providecommand{\textding}[1]{#1}
\providecommand{\Oij}{\(\mathrm{O};\vec{\imath},\vec{\jmath}\)}
\providecommand{\Oijk}{\(\mathrm{O};\vec{\imath},\vec{\jmath},\vec{k}\)}
\providecommand{\Ouv}{\(\mathrm{O};\vec{u},\vec{v}\)}
\providecommand{\intervFF}[2]{\ensuremath{[#1;#2]}}
\renewcommand{\euro}{\ifmmode\text{\texteuro}\else\texteuro\fi}
\DeclareUnicodeCharacter{00B0}{\ifmmode^{\circ}\else\textdegree\fi}
\DeclareUnicodeCharacter{20AC}{\ifmmode\text{\texteuro}\else\texteuro\fi}
% Énoncés : Éducation nationale. Transcription et corrigés : APMEP, Denis Vergès et contributeurs. Découpage et classement : CoopMaths, Sébastien Lozano, Éric Elter et contributeurs. Adaptation de la présentation : Maths Explorer.
% CC BY-SA : https://coopmaths.fr/www/about/
% https://coopmaths.fr/alea/static/bac/2023/tex/bac_2023_03_sujet1_polynesie_1.tex
\begin{document}
\noindent\textit{APMEP / CoopMaths. CC BY-SA. Présentation adaptée par Maths Explorer.}\par\medskip
\medskip

\emph{Les parties {\rm A} et {\rm B} peuvent être traitées indépendamment}

\medskip

Les utilisateurs de vélo d'une ville sont classés en deux catégories disjointes:

\begin{itemize}
\item[$\bullet~~$]ceux qui utilisent le vélo dans leurs déplacements professionnels ;
\item[$\bullet~~$]ceux qui utilisent le vélo uniquement pour leurs loisirs.
\end{itemize}

\medskip

Un sondage donne les résultats suivants:

\begin{itemize}
\item[$\bullet~~$]21\,\% des utilisateurs ont moins de 35 ans.

Parmi eux, 68\,\% utilisent leur vélo uniquement pour leurs loisirs alors que les autres l'utilisent dans leurs déplacements professionnels ;
\item[$\bullet~~$]parmi les 35 ans ou plus, seuls 20\,\% utilisent leur vélo dans leurs déplacements professionnels, les autres l'utilisent uniquement pour leurs loisirs.
\end{itemize}

On interroge au hasard un utilisateur de vélo de cette ville.

Dans tout l'exercice on considère les évènements suivants:

\begin{itemize}
\item[$\bullet~~$]$J$ : \og la personne interrogée a moins de 35 ans \fg{} ;
\item[$\bullet~~$]$T$ : \og la personne interrogée utilise le vélo dans ses déplacements professionnels \fg{} ;
\item[$\bullet~~$]$\overline{J}$ et $\overline{T}$ sont les évènements contraires de $J$ et $T$.
\end{itemize}

\bigskip

\textbf{Partie A}

\medskip

\begin{enumerate}
\item Calculer la probabilité que la personne interrogée ait moins de $35$~ans et utilise son vélo dans ses déplacements professionnels.

\emph{On pourra s'appuyer sur un arbre pondéré.}
\item Calculer la valeur exacte de la probabilité de $T$.
\item On considère à présent un habitant qui utilise son vélo dans ses déplacements professionnels.

Démontrer que la probabilité qu'il ait moins de 35 ans est $0,30$ à $10^{-2}$ près.
\end{enumerate}

\bigskip

\textbf{Partie B}

\medskip

Dans cette partie, on s'intéresse uniquement aux personnes utilisant leur vélo dans leurs déplacements professionnels.

\medskip

On admet que 30\,\% d'entre elles ont moins de $35$~ans.

On sélectionne au hasard parmi elles un échantillon de $120$~personnes auxquelles on va soumettre un questionnaire supplémentaire. 

On assimile la sélection de cet échantillon à un tirage aléatoire avec remise.

\medskip

On demande à chaque individu de cet échantillon son âge.

\medskip

X représente le nombre de personnes de l'échantillon ayant moins de $35$~ans.

\medskip

\emph{Dans cette partie, les résultats seront arrondis à $10^{-3}$ près.}

\medskip

\begin{enumerate}
\item Déterminer la nature et les paramètres de la loi de probabilité suivie par $X$.
\item Calculer la probabilité qu'au moins $50$ utilisateurs de vélo parmi les $120$ aient moins de $35$~ans.
\end{enumerate}

\bigskip

\end{document}
